%-------------------------
% Resume in LaTeX (Japanese)
% Author : Kenta Utsunomiya
% Based on: https://www.overleaf.com/articles/vidushi-wahals-cv/wpcddrydqwsj
% License : MIT
%------------------------

\documentclass[11pt]{article}

\usepackage{luatexja-fontspec}
\setmainjfont{Noto Sans CJK JP}

\usepackage{geometry}
\geometry{letterpaper, margin=0.75in}

\usepackage{titlesec}
\usepackage[usenames,dvipsnames]{color}
\usepackage{verbatim}
\usepackage{enumitem}
\usepackage[hidelinks]{hyperref}
\usepackage{fancyhdr}

\pagestyle{fancy}
\fancyhf{} % clear all header and footer fields
\fancyfoot{}
\renewcommand{\headrulewidth}{0pt}
\renewcommand{\footrulewidth}{0pt}

\urlstyle{same}

\raggedbottom
\raggedright
\setlength{\tabcolsep}{0in}

% Sections formatting
\titleformat{\section}{
  \vspace{-8pt}\raggedright\large\bfseries
}{}{0em}{}[\color{black}\titlerule \vspace{-4pt}]

%-------------------------
% Custom commands
\newcommand{\resumeItem}[1]{
  \item\small{
    {#1 \vspace{-1pt}}
  }
}

\newcommand{\resumeSubheading}[4]{
  \item
    \begin{tabular*}{1.0\textwidth}[t]{l@{\extracolsep{\fill}}r}
      \textbf{#1} & #2 \\
      \textit{\small#3} & \textit{\small #4} \\
    \end{tabular*}\vspace{-2pt}
}

\newcommand{\resumeSubItem}[2]{\resumeItem{\textbf{#1}{: #2}}\vspace{-2pt}}

\newcommand{\resumeCert}[2]{
  \item
    \begin{tabular*}{1.0\textwidth}[t]{l@{\extracolsep{\fill}}r}
      \textbf{#1} & \textit{\small #2} \\
    \end{tabular*}\vspace{-2pt}
}

\renewcommand{\labelitemii}{$\circ$}

\newcommand{\resumeSubHeadingListStart}{\begin{itemize}[leftmargin=0.0in, label={}]}
\newcommand{\resumeSubHeadingListEnd}{\end{itemize}}
\newcommand{\resumeItemListStart}{\begin{itemize}[leftmargin=1.5em, noitemsep, topsep=1pt]}
\newcommand{\resumeItemListEnd}{\end{itemize}\vspace{-2pt}}

%-------------------------------------------
%%%%%%  CV STARTS HERE  %%%%%%%%%%%%%%%%%%%%%%%%%%%%

\begin{document}

%----------HEADING-----------------
\begin{center}
  \textbf{\huge Kenta Utsunomiya (宇都宮 健太)} \\ \vspace{3pt}
  \small ダブリン、アイルランド \ $|$ \
  \href{mailto:utsu.k.k@gmail.com}{\underline{utsu.k.k@gmail.com}} \ $|$ \
  \href{https://ayashiro.github.io/}{\underline{ayashiro.github.io}} \ $|$ \
  \href{https://www.linkedin.com/in/kenta-utsunomiya-301411153/}{\underline{linkedin.com/in/kenta-utsunomiya}} \ $|$ \
  \href{https://github.com/ayashiro}{github.com/ayashiro}
\end{center}
\vspace{-8pt}

%-----------EXPERIENCE-----------------
\section{職務経歴}
  \resumeSubHeadingListStart

    \resumeSubheading
      {Google Ireland}{ダブリン、アイルランド}
      {ソフトウェアエンジニア (L4)}{2021年3月 -- 現在}
      \resumeItemListStart
        \resumeItem{財務データ移行プロジェクトのテックリードとして、Cloud Composer (Airflow) を用いたパイプラインを設計し、分散データセット月間1PBをBigQueryに集約。エンジニア3名・アナリスト4名・PM1名から成る8名の部門横断チームを統率した。}
        \resumeItem{全社的な株式・ストックグラント管理向け社内Webアプリケーションの開発をリード。技術設計ドキュメントを作成し、事業側ステークホルダーとの合意形成を主導した。}
        \resumeItem{四半期ごとの財務レポーティングETLパイプラインを100以上のデータソースにわたり自動化された月次ワークフローへ再構築し、処理時間を1週間から1日に短縮した。}
        \resumeItem{Googleのエンジニア採用面接官として100名以上の候補者を評価。コーディング面接問題委員会のコアメンバーとして、競技プログラミング形式のアルゴリズム問題を設計(20\%プロジェクト)。}
      \resumeItemListEnd

    \resumeSubheading
      {Indeed Japan / Recruit Holdings}{東京、日本 \& オースティン、テキサス}
      {ソフトウェアエンジニア (Lv1 $\rightarrow$ Lv2)}{2016年4月 -- 2021年3月}
      \resumeItemListStart
        \resumeItem{Java製のミッションクリティカルなP2Pファイル配信ミドルウェアを運用・スケーリングし、求人検索インデックス(Jobindex)を日次5TB以上、世界各地のデータセンター間で転送。圧縮処理と独自にフォークしたBitTorrentクライアントによりダウンロードデータ量を削減し、データセンター間の帯域使用量を80\%削減。人事評価で上位5\%を獲得した。}
        \resumeItem{Jaeger・LightstepなどのDistributed Tracingソリューションをベンチマーク・評価し、Kubernetesクラスタ上にPrometheus/Grafanaと共に導入。全社的なマイクロサービスの可観測性(オブザーバビリティ)を向上させた。}
        \resumeItem{テキサス州オースティンで開催されたIndeed University インキュベータープログラムに選抜。営業部門・法務部門と連携し、前科を持つ求職者を支援するサービス「Indeed Restart」の企画から開発・ローンチまでを3ヶ月で主導した。}
        \resumeItem{複数の新卒エンジニアおよびインターンのメンタリングを担当し、厳格なコードレビュー体制とDockerベースのCIテストワークフローを整備した。}
      \resumeItemListEnd

  \resumeSubHeadingListEnd

%-----------EDUCATION-----------------
\section{学歴}
  \resumeSubHeadingListStart
    \resumeSubheading
      {東京工業大学}{東京、日本}
      {数理・計算科学専攻 修士課程 修了 \textnormal{(修士論文: 数理物理学)}}{2013年10月 -- 2016年9月}
    \resumeSubheading
      {東京工業大学}{東京、日本}
      {情報科学科 学士課程 卒業 \textnormal{(研究テーマ: 複雑ネットワークとグラフマイニング)}}{2009年4月 -- 2013年9月}
  \resumeSubHeadingListEnd

%--------SKILLS & TECHNOLOGIES------------
\section{スキル・技術}
 \resumeSubHeadingListStart
   \item{
     \textbf{プログラミング言語}{: Java, C++, Python, Go, SQL, MATLAB}
   }
   \item{
     \textbf{クラウド・インフラ}{: Google Cloud Platform (GCP), BigQuery, Cloud Composer (Airflow), Kubernetes, Docker, Terraform}
   }
   \item{
     \textbf{分散システム・バックエンド}{: 分散システム, P2Pネットワーク, マイクロサービス, Prometheus, Grafana, Linux, Git}
   }
   \item{
     \textbf{言語}{: 日本語(母語), 英語(ビジネスレベル / IELTS 7.5, TOEIC 900点), 中国語(HSK5級)}
   }
 \resumeSubHeadingListEnd

%--------CERTIFICATIONS & HONORS------------
\section{資格・受賞}
  \resumeSubHeadingListStart
    \resumeCert{統計検定 1級}{2019年11月}
    \resumeCert{実用数学技能検定 1級}{2012年12月}
  \resumeSubHeadingListEnd

%-------------------------------------------
\end{document}
